% Generated draft: requires mathematical and source review.
\documentclass[11pt]{article}
\usepackage[margin=25mm]{geometry}
\usepackage{fontspec}
\setmainfont{Latin Modern Roman}
\usepackage{amsmath,amssymb,mathtools,unicode-math}
\setmathfont{Latin Modern Math}
\usepackage{enumitem,needspace,xurl,hyperref}
\setlength{\parindent}{0pt}
\setlength{\parskip}{5pt}
\setlength{\emergencystretch}{3em}
\newcommand{\N}{\mathbb{N}}
\newcommand{\Z}{\mathbb{Z}}
\newcommand{\Q}{\mathbb{Q}}
\newcommand{\R}{\mathbb{R}}
\newcommand{\C}{\mathbb{C}}
\newcommand{\F}{\mathbb{F}}
\newcommand{\A}{\mathbb{A}}
\newcommand{\PP}{\mathbb{P}}
\newcommand{\E}{\mathbb{E}}
\newcommand{\eps}{\varepsilon}
\newcommand{\dd}{\,\mathrm{d}}
\begin{document}
\section*{UnsolvedMath ID: 20002338}
\textbf{Source problem code:} catalogue entry
\subsection{Definitions and mathematical statement}
For every prime $p$ and every integer $x$, Fermat's little theorem states
\[
 x^p\equiv x\pmod p,
\]
or equivalently $a^{p-1}\equiv1\pmod p$ when $p\nmid a$. The source asks for proofs of this theorem and a ``simple'' definition of $n!\bmod p$, allowing exponentiation. It sketches three routes: cyclicity of $\F_p^\times$, using the bound $|\{x\in\F_p:P(x)=0\}|\le\deg P$ for nonzero polynomials $P$; the permutation of $\{1,\ldots,p-1\}$ induced by multiplication by a nonzero residue $a$, giving
\[
 (p-1)!\equiv\prod_{i=1}^{p-1}(ai)
 \equiv a^{p-1}(p-1)!\pmod p;
\]
and the identity $(x+y)^p\equiv x^p+y^p\pmod p$. Cancellation in the product argument uses $p\nmid(p-1)!$. The supplied research aid reports completed versions of these arguments and further proofs using Lagrange's theorem and the orbits of length-$p$ words under cyclic rotation.

For the factorial-definition request, the aid adopts standard $\N$ with addition, multiplication, order and the binary exponentiation function. The following makes its bounded single-code formulation explicit. For $p\ge2$, define the base-$p$ digit relation by
\[
 \operatorname{Digit}_p(C,j,u)\ \Longleftrightarrow\
 u<p\ \wedge\ \exists a\le C\ \exists b<p^j\
 \bigl(C=ap^{j+1}+up^j+b\bigr).
\]
Then, for all $n\in\N$, $p\ge2$ and $r\in\N$, the reported definition of $r=n!\bmod p$ is
\[
\begin{split}
 r<p\ \wedge\ \exists C<p^{n+1}\ \bigl(&\operatorname{Digit}_p(C,0,1)
 \wedge\operatorname{Digit}_p(C,n,r)\\
 &\wedge\ \forall j<n\ \exists u<p\ \exists v<p\ \exists q\le j\
 [\operatorname{Digit}_p(C,j,u)\wedge
 \operatorname{Digit}_p(C,j+1,v)\\
 &\hspace{42mm}\wedge\ (j+1)u=v+qp]\bigr).
\end{split}
\]
Every variable here ranges over $\N$. The code stores the residues of $0!,1!,\ldots,n!$ as its successive digits, and the quotient bound is $q\le j$. The assertion is that this bounded formula is correct and functional in standard arithmetic for every modulus $p\ge2$, including composite moduli, with the boundary case $n=0$ giving $r=1$. Primality is needed for Fermat's theorem, not for this definition. The source does not specify a weak formal theory in which correctness must be proved, a permitted quantifier class without exponentiation, or a precise complexity criterion for ``simple.'' The bounded formula is one explicit interpretation of that request, not a claim of an existential polynomial definition or of provability in every weak arithmetic.
\subsection{Short English statement}
Give proofs of Fermat's little theorem and a simple definition of factorial residues. With exponentiation available, the research aid uses one bounded base-$p$ code for successive factorial residues; the definition also works for composite moduli.
\subsection{Sources}
\begin{itemize}
\item Paola D'Aquino, Question 1 in the American Institute of Mathematics workshop \emph{Extensions of Hilbert's Tenth Problem}, 21--25 March 2005, page 1 as identified in the supplied verification, \url{https://aimath.org/WWN/hilberts10th/hilberts10th.pdf}.
\item Local UnsolvedMath registry, permanent ID 20002338, source code \texttt{catalogue entry}, canonical location \texttt{aim-logic-notes.json}, \texttt{notes[113]}, record dated 14 August 2026. Its verification reports lost exponents and damaged group and product notation in the extraction; the formulas here restore the meanings expressly described there. The original literature field is empty.
\item Supplied research attempt 1: partial result, problem status at run solved, rigor labelled proved, novelty confidence low. The single-code definition and additional proofs are machine-generated research aid, not peer reviewed, and require independent verification before citation; no separate bibliographic source is supplied for them.
\end{itemize}
\end{document}
